\section{Partial Derivatives}
\label{sec:partial-derivatives}

\subsection{Partial Derivatives}

Many physical quantities depend on several variables.

Suppose

\[
f(x,y,z)=x^2+y^2+z.
\]

Differentiating only with respect to

\[
x,
\]

while treating

\[
y,z
\]

as constants gives

\[
\boxed{
	\frac{\partial f}{\partial x}=2x.
}
\]

Similarly,

\[
\frac{\partial f}{\partial y}=2y,
\]

and

\[
\frac{\partial f}{\partial z}=1.
\]

Partial derivatives describe how a function changes along one coordinate
direction while all other coordinates remain fixed.

%%%%%%%%%%%%%%%%%%%%%%%%%%%%%%%%%%%%%%%%%%%%%%%%%%%%%%%%%%%%%%%%%%%%%%%%%%%%%%%%

\subsection{Total Derivative}

Suppose

\[
f=f(x(t),y(t)).
\]

Now every variable depends upon time.

The total derivative becomes

\[
\boxed{
	\frac{df}{dt}
	=
	\frac{\partial f}{\partial x}
	\frac{dx}{dt}
	+
	\frac{\partial f}{\partial y}
	\frac{dy}{dt}.
}
\]

This is the multivariable chain rule.

%%%%%%%%%%%%%%%%%%%%%%%%%%%%%%%%%%%%%%%%%%%%%%%%%%%%%%%%%%%%%%%%%%%%%%%%%%%%%%%%

\subsection{The Differential Operator}

The symbol

\[
\boxed{
	\nabla
}
\]

is called

\begin{itemize}
	\item the del operator,
	\item the nabla operator,
	\item the vector differential operator.
\end{itemize}

It is defined as

\[
\boxed{
	\nabla
	=
	\hat{\mathbf x}\frac{\partial}{\partial x}
	+
	\hat{\mathbf y}\frac{\partial}{\partial y}
	+
	\hat{\mathbf z}\frac{\partial}{\partial z}.
}
\]

Unlike an ordinary vector,

\[
\nabla
\]

contains differentiation operators instead of numerical components.

%%%%%%%%%%%%%%%%%%%%%%%%%%%%%%%%%%%%%%%%%%%%%%%%%%%%%%%%%%%%%%%%%%%%%%%%%%%%%%%%
